\documentclass[11pt, a4paper]{article}

\usepackage[T1]{fontenc}
\usepackage[utf8]{inputenc}
\usepackage{tgpagella}
\usepackage[spanish, es-noshorthands]{babel}
\usepackage{amsmath, amssymb, bm}
\usepackage{tabularx}
\usepackage{booktabs}
\usepackage[most]{tcolorbox}
\usepackage[margin=2.5cm]{geometry}
\usepackage{ragged2e}
\usepackage{fancyhdr}

\pagestyle{fancy}
\fancyhf{}
\setlength{\headheight}{16pt}
\lhead{\textbf{UCB Sede Tarija} | Ing. Mecatrónica}
\chead{}
\rhead{IMT-121 Circuitos Electrónicos I | Auditoría de Progreso}
\lfoot{Prof: M.Sc. Ing. Bernardo Quiroga Turdera}
\rfoot{Página \thepage}
\renewcommand{\headrulewidth}{0.4pt}

\definecolor{ucbblue}{RGB}{0, 51, 102}

\begin{document}

\begin{tcolorbox}[colback=ucbblue!5!white, colframe=ucbblue, title=\textbf{Auditoría Docente de Evidencia Analítica RC (Readiness Gate)}]
    \centering \textbf{Decisión de Transición hacia Entorno de Simulación}
\end{tcolorbox}

\noindent \textit{Instrucciones: Evalúe el desempeño oral/escrito del estudiante basado en los ejercicios del repositorio W11S1. Marque para identificar la raíz de la falla.}

\vspace{0.4cm}
\renewcommand{\arraystretch}{1.6}
\noindent
\begin{tabularx}{\textwidth}{>{\RaggedRight\arraybackslash}X >{\centering\arraybackslash}p{2.5cm} >{\centering\arraybackslash}p{2cm} >{\centering\arraybackslash}p{2.5cm}}
\toprule
\textbf{Criterio Estructural Requerido} & \textbf{Demostrado} & \textbf{Parcial} & \textbf{No Logrado} \\ 
\midrule
Mantiene y justifica la continuidad de voltaje $v_C(0^+) = v_C(0^-)$ & $\square$ & $\square$ & $\square$ \\ \addlinespace
Identifica la red en $t \to \infty$ asumiendo la rama del capacitor abierta & $\square$ & $\square$ & $\square$ \\ \addlinespace
Obtiene $R_{eq}$ aislando el puerto y apagando fuentes energéticas & $\square$ & $\square$ & $\square$ \\ \addlinespace
Arma la ecuación exponencial respetando dimensionalidad en $\tau$ & $\square$ & $\square$ & $\square$ \\ \addlinespace
Verifica matemáticamente los límites $t=0^+$ y $t \to \infty$ en su fórmula & $\square$ & $\square$ & $\square$ \\ \addlinespace
Distingue el tiempo absoluto ($t$) del tiempo relativo a la conmutación ($s$) & $\square$ & $\square$ & $\square$ \\ 
\bottomrule
\end{tabularx}

\vspace{0.8cm}
\begin{tcolorbox}[colback=gray!5, colframe=gray!50, title=\textbf{Toma de Decisión Docente}]
En base al diagnóstico superior, categorice el estado del estudiante o del grupo:
\begin{itemize}
    \item[$\square$] \textbf{\texttt{READY\_FOR\_TRANSIENT\_SIMULATION}}: Errores menores o aritméticos que no afectan la comprensión física del RC de primer orden. Autorizado para proceder a verificación \texttt{.tran}.
    \item[$\square$] \textbf{\texttt{NEEDS\_RC\_REINFORCEMENT}}: Confusiones graves persistentes en principios de inercia eléctrica, KVL o $R_{eq}$. Permanecer en lápiz y papel.
\end{itemize}
\end{tcolorbox}

\end{document}
